\section{Robustness and limitations}
\label{sec:robust}

\paragraph{Contract type and sub-period.} Table~\ref{tab:h3h4} repeats the strategies by contract
type (H3) and sub-period (H4). The contract split is the important one (Section~\ref{sec:review}).
The sub-period split cannot test H4 as intended: 2026 alone covers under a year of each chip's
life, and before 2023 only three chips are on sale, too few clusters to report.

\begin{table}[htbp]
\centering
\caption{Age profile by contract type and sub-period}
\label{tab:h3h4}
\begin{adjustbox}{max width=\textwidth}
% Generated by gel.viz.hedonic_figures from p5_summary.csv. Do not edit.
% Rates in log points per year; early - late = rate(0-2) - rate(2-4); H1 predicts < 0.
\begin{tabular}{llccc}
\hline
Sample & Strategy & rate, yrs 0--2 & rate, yrs 2--4 & early $-$ late (H1: $<0$) \\
\hline
On demand & 1a quality index (TFLOP/s) & $-$0.07 [$-$0.47, 0.34] & 0.10 [$-$0.16, 0.37] & $-$0.17 [$-$0.70, 0.37] \\
On demand & 2 frontier-relative & $-$0.21 [$-$0.48, 0.07] & $-$0.09 [$-$0.36, 0.17] & $-$0.12 [$-$0.33, 0.08] \\
On demand & 4 shape only & -- & -- & $-$0.13 [$-$0.40, 0.13] \\
Reserved, 1 year & 1a quality index (TFLOP/s) & 0.81 [$-$0.54, 2.19] & 0.18 [$-$0.08, 0.45] & 0.63 [$-$0.78, 2.03] \\
Reserved, 1 year & 2 frontier-relative & 0.08 [$-$0.18, 0.34] & $-$0.06 [$-$0.23, 0.11] & 0.14 [$-$0.23, 0.50] \\
Reserved, 1 year & 4 shape only & -- & -- & 0.76 [$-$1.15, 2.66] \\
Reserved, 3 years & 1a quality index (TFLOP/s) & 0.19 [$-$0.07, 0.45] & 0.10 [$-$0.18, 0.37] & 0.09 [$-$0.05, 0.24] \\
Reserved, 3 years & 2 frontier-relative & $-$0.03 [$-$0.07, 0.01] & $-$0.12 [$-$0.28, 0.04] & 0.09 [$-$0.13, 0.31] \\
Reserved, 3 years & 4 shape only & -- & -- & $-$0.20 [$-$0.80, 0.38] \\
Spot & 1a quality index (TFLOP/s) & 0.46 [0.11, 0.81] & $-$0.19 [$-$0.32, $-$0.04] & 0.65 [0.19, 1.10] \\
Spot & 2 frontier-relative & 0.50 [0.18, 0.83] & $-$0.26 [$-$0.42, $-$0.09] & 0.76 [0.26, 1.26] \\
Spot & 4 shape only & -- & -- & 0.87 [0.34, 1.41] \\
Without Google, July 2026 & 1a quality index (TFLOP/s) & 0.28 [$-$0.21, 0.76] & $-$0.03 [$-$0.24, 0.18] & 0.31 [$-$0.50, 1.09] \\
Without Google, July 2026 & 2 frontier-relative & 0.11 [$-$0.12, 0.35] & $-$0.15 [$-$0.39, 0.10] & 0.26 [0.05, 0.47] \\
Without Google, July 2026 & 4 shape only & -- & -- & 0.47 [0.36, 0.57] \\
2023--2025 & 1a quality index (TFLOP/s) & 0.26 [$-$0.11, 0.63] & $-$0.27 [$-$0.43, $-$0.12] & 0.54 [0.19, 0.88] \\
2023--2025 & 2 frontier-relative & 0.11 [$-$0.09, 0.32]$^\dagger$ & $-$0.37 [$-$1.62, 0.86]$^\dagger$ & 0.48 [0.34, 0.63] \\
2023--2025 & 4 shape only & -- & -- & 0.48 [0.28, 0.68] \\
2026 & 1a quality index (TFLOP/s) & $-$0.36 [$-$53.25, 53.09] & $-$0.09 [$-$0.29, 0.12] & $-$0.27 [$-$52.97, 52.98] \\
2026 & 2 frontier-relative & $-$6.60 [$-$47.69, 35.77]$^\dagger$ & 0.82 [$-$2.86, 4.12]$^\dagger$ & $-$7.42 [$-$44.76, 30.91] \\
2026 & 4 shape only & -- & -- & $-$0.61 [$-$11.87, 10.71] \\
Provider-chip-months weighted equally & 1a quality index (TFLOP/s) & 0.21 [$-$0.21, 0.62] & $-$0.03 [$-$0.18, 0.11] & 0.24 [$-$0.15, 0.64] \\
Provider-chip-months weighted equally & 2 frontier-relative & 0.10 [$-$0.22, 0.43] & $-$0.08 [$-$0.17, 0.01] & 0.19 [$-$0.11, 0.47] \\
Provider-chip-months weighted equally & 4 shape only & -- & -- & 0.29 [0.07, 0.50] \\
\hline
\end{tabular}
% dagger: within-chip age span under 4 years with chip effects (extrapolated).

\end{adjustbox}
\par\smallskip\footnotesize As Table~\ref{tab:strategies}. $\dagger$: within-chip age span under
four years with chip effects (extrapolated). Sub-samples with fewer than six chips are not reported.
\end{table}

\paragraph{Weighting and composition.} Google Cloud is \nGoogleShare{} of the estimation sample.
Weighting each provider equally lowers the pooled shape to \nShapeEqualProv, and dropping Google
leaves \nShapeNoGoogle. Strategy 1's early-years rate depends on the H200: dropping it moves the
rate from \nOneARateEarly{} to \nOneARateEarlyNoHTwoHundred.

\paragraph{Posted versus transacted prices.} The panel is posted prices, overwhelmingly hyperscaler
list prices, and \nListUnchanged{} of on-demand and reserved month-to-month transitions are
unchanged. Large buyers pay negotiated prices that are not observed. \ajcheck{The results are
therefore about the posted-price market; a transacted-price panel, such as the marketplace data of
\citet{huang2026gpu}, would be the natural complement.}

\paragraph{Age zero.} Age is measured from Epoch's release month, which precedes the first cloud
listing by several months for some chips. The event study's sensitivity to that dating is reported
in Section~\ref{sec:events}.

\paragraph{Operating cost.} Hosting, network and staff costs are excluded for lack of a sourced
figure, so values are upper bounds and $T^*$ is biased upward.
